\ExerciseSection

* 1) Let \(\mathbf{I}\) be the closed cube of \(\mathbf{R}^n\) which is the product of \(n\) intervals identical with \([-\frac{1}{2}\pi, \frac{1}{2}\pi]\) . To each point \(x = (x_i)\) of \(\mathbf{I}\) we make correspond the point \(y = (y_j)\) of \(\mathbf{R}^{n+1}\) such that

\[
y _ {1} = \sin x _ {1}
\]

\[
{y _ {2}} {= \cos x _ {1} \sin x _ {2}}
\]

\[
y _ {p} = \cos x _ {1} \cos x _ {2} \dots \cos x _ {p - 1} \sin x _ {p} (2 \leqslant p \leqslant n - 1)
\]

\[
y _ {n} = \cos x _ {1} \cos x _ {2} \dots \cos x _ {n - 1} \sin 2 x _ {n}
\]

\[
y _ {n + 1} = \cos x _ {1} \cos x _ {2} \dots \cos x _ {n - 1} \cos 2 x _ {n}.
\]

Show that the image of I under this mapping is the sphere \(S_{n}\) , and that the restriction of this mapping to the interior of I is a homeomorphism onto a dense open set in \(S_{n}\) . *

\begin{enumerate}
\def\labelenumi{\arabic{enumi})}
\setcounter{enumi}{1}
\tightlist
\item
  Define a homeomorphism of \(\mathbf{S}_p \times \mathbf{S}_q\) onto a subset of \(\mathbf{S}_{p + q + 1}\) {[}note that the equation of \(\mathbf{S}_{p + q + 1}\) is
\end{enumerate}

\[
\left. \left(x _ {1} ^ {2} + \dots + x _ {p + 1} ^ {2}\right) + \left(x _ {p + 2} ^ {2} + \dots + x _ {p + q + 2} ^ {2}\right) = \mathrm{I} \right].
\]

\begin{enumerate}
\def\labelenumi{\arabic{enumi})}
\setcounter{enumi}{2}
\tightlist
\item
  \begin{enumerate}
  \def\labelenumii{\alph{enumii})}
  \tightlist
  \item
    Show that, if \(f\) is a continuous mapping of \(\mathbf{S}_1\) into \(\mathbf{R}^n\) , \(f\) can be extended to a continuous mapping of \(\mathbf{B}_2\) into \(\mathbf{R}^n\) {[}map the point \(tx \in \mathbf{B}_2\) ( \(t \geqslant 0\) , \(x \in S_1\) ) to \(tf(x)\) {]}.
  \end{enumerate}
\end{enumerate}

\begin{enumerate}
\def\labelenumi{\alph{enumi})}
\setcounter{enumi}{1}
\tightlist
\item
  Deduce that, if \(f\) is a continuous mapping of \(\mathbf{S}_1\) into \(\mathbf{S}_n\) such that \(f(\mathbf{S}_1) \neq \mathbf{S}_n\) , then \(f\) can be extended to a continuous mapping of \(\mathbf{B}_2\) into \(\mathbf{S}_n\) {[}use a stereographic projection whose vertex does not lie in \(f(\mathbf{S}_1)\) {]}. (*)
\end{enumerate}

\begin{enumerate}
\def\labelenumi{\arabic{enumi})}
\setcounter{enumi}{3}
\tightlist
\item
  Show that there is no homeomorphism of \(S_{1}\) into R. (Observe that the complement in \(S_{1}\) of any point of \(S_{1}\) is connected, and that every connected subset of R is an interval.)
\end{enumerate}

Deduce that a homeomorphism of \(S_{1}\) onto a subspace of \(S_{1}\) is necessarily a homeomorphism of \(S_{1}\) onto \(S_{1}\) (use Proposition 4 of no. 4).

\begin{enumerate}
\def\labelenumi{\arabic{enumi})}
\setcounter{enumi}{4}
\item
  Show that, if \(n > 1\) , the sphere \(\mathbf{S}_n\) is not homeomorphic to the circle \(\mathbf{S}_1\) (cf.~§ 1, Exercise 8).
\item
  Identifying \(S_{1}\) with the torus T, the quotient of R by the relation \(x \equiv y \pmod{1}\) (no. 4, Proposition 4, Corollary 2), let \(\varphi\) denote the canonical mapping of R onto \(S_{1}\) . A continuous mapping f of a topological space E into \(S_{1}\) is inessential if there exists a continuous mapping g of E into R such that \(f = \varphi \circ g\) . (**) A mapping which is not inessential is said to be essential.
\end{enumerate}

Show that the identity mapping of \(S_{1}\) onto \(S_{1}\) is essential (use Exercise 4).

\begin{enumerate}
\def\labelenumi{\arabic{enumi})}
\setcounter{enumi}{6}
\item
  Show that there exists an entourage U of the uniformity of \(S_{1}\) such that, if f is an inessential mapping of a topological space E into \(S_{1}\) , every continuous mapping \(f' : E \to S_{1}\) such that \((f(x), f'(x)) \in U\) for all \(x \in E\) is also inessential.
\item
  Show that there exists no continuous mapping \(f\) of \(\mathbf{B}_2\) onto \(\mathbf{S}_1\) which extends the identity mapping of \(\mathbf{S}_1\) (using Exercise 7, show that \(f\) , restricted to the circle with centre o and radius \(r \leqslant i\) , would always be an inessential mapping of this circle into \(S_{1}\) , and hence obtain a contradiction when \(r = i\) ) .
\item
  Let E be a topological space containing more than one point. Show that there exists a continuous mapping f of \(S_{1}\) into \(F = E \times S_{1}\) such that \(f(S_{1}) \neq F\) and which does not extend to a continuous mapping of \(B_{2}\) into F (use Exercise 8). Deduce that, if n \textgreater{} 1, \(S_{n}\) , \(R^{n}\) and \(B_{n}\) are not homeomorphic to any space of the form \(E \times S_{1}\) , where E is any topological space (see Exercise 3). In particular, if n \textgreater{} 1, \(S_{n}\) is not homeomorphic to \(S_{1}^{n}\) , and for all n, \(B_{n}\) is not homeomorphic to \(S_{1}^{n}\) .
\end{enumerate}

Show likewise that \(R^{2}\) is not homeomorphic to the complement of a point in \(R^{2}\) , and that \(S_{2}\) is not homeomorphic to \(B_{2}\) (if false, \(R^{2}\) would be homeomorphic to the product of \(S_{1}\) and the interval \([0, 1]\) ).

\begin{enumerate}
\def\labelenumi{\arabic{enumi})}
\setcounter{enumi}{9}
\tightlist
\item
  Let \(H_{n,p,q}\) denote the ``quadric'' in \(R^{n}\) defined by the equation
\end{enumerate}

\[
x _ {1} ^ {2} + x _ {2} ^ {2} + \dots + x _ {p} ^ {2} - x _ {p + 1} ^ {2} - \dots - x _ {p + q} ^ {2} = \mathrm{I} \quad (p + q \leqslant n).
\]

Show that \(\mathbf{H}_{n,p,q}\) is homeomorphic to \(\mathbf{S}_{p - 1}\times \mathbf{R}^{n - p}\) .

\begin{enumerate}
\def\labelenumi{\arabic{enumi})}
\setcounter{enumi}{10}
\tightlist
\item
  Let \(C_{n,p}\) be the ``quadric cone'' in \(R^{n}\) defined by the equation
\end{enumerate}

\[
x _ {1} ^ {2} + x _ {2} ^ {2} + \dots + x _ {p} ^ {2} - x _ {p + 1} ^ {2} - \dots - x _ {n} ^ {2} = 0 \quad (\mathrm{I} \leqslant p \leqslant n - \mathrm{I}).
\]

Show that the complement of \(\{0\}\) in \(\mathbf{C}_{n,p}\) is homeomorphic to \(\mathbf{S}_{p-1} \times \mathbf{S}_{n-p-1} \times \mathbf{R}\) .

* 12) A subset E of \(\mathbf{R}^n\) which contains the origin o is said to be starlike (with respect to o) if, for all \(x \in \mathbf{E}\) and all \(t \in [0, 1]\) , we have \(tx \in \mathbf{E}\) . The intersection of E with a closed ray of origin o is either the whole ray or a segment of which o is one end-point. The shell of E is the set K consisting of the non-zero end-points of these segments, together with o if there exists a ray whose intersection with E consists of o alone. In what follows we shall assume that \(o \notin K\) .

\begin{enumerate}
\def\labelenumi{\alph{enumi})}
\item
  Show that the shell of E is contained in the frontier of E. Give an example in which these two sets are different.
\item
  Show that, if the shell K of E is compact, there is a homeomorphism of \(R^{n}\) onto itself which maps K onto \(S_{n-1}\) , \(\overline{E}\) onto \(B_{n}\) and the interior of E onto the interior of \(B_{n}\) (map each point \(x \in K\) to its central projection on \(S_{n-1}\) , and then extend this mapping to the whole of \(R^{n}\) ). Deduce that the frontier of E coincides with its shell, and that the interior of E is the set of points tx, where \(x \in K\) and \(t \in [0, I[\) .
\item
  If E is unbounded and if its frontier coincides with its shell, then the interior of E is homeomorphic to \(R^{n}\) , and its shell K is homeomorphic to an open subset of \(S_{n-1}\) {[}show that the image of E under the homeomorphism \(x \to x/(1 + ||x||)\) satisfies the conditions of b){]}.
\item
  Give an example of an unbounded starlike set E whose shell is closed but is not identical with the frontier of E.
\end{enumerate}

\begin{enumerate}
\def\labelenumi{\arabic{enumi})}
\setcounter{enumi}{12}
\tightlist
\item
  Show that, in the space \(S_{n}\) , the frontier of a non-empty open set whose exterior is non-empty has the power of the continuum (use Proposition 4 of no. 4).
\end{enumerate}

